%------------------------------------------------------------------------------%
\documentclass[fleqn,utf8,aspectratio=169,12pt]{beamer}

\usepackage{integracao_e_series}

\title{Séries Numéricas -- Teste da Razão}

%------------------------------------------------------------------------------%
\begin{document}

\addtitle

%------------------------------------------------------------------------------%
\section{Teste da Razão}

%------------------------------------------------------------------------------%
\begin{proofframe}{Motivação -- Analise da Série Geométrica}
	Considere a série geométrica com $a$ e $r$ positivos
	\[
    \sum_{n=1}^{\infty} ar^{n-1}
  \]
	\pause
	Podemos comparar o crescimento de seus termos calculando a razão
	\[
    \rho
    = \frac{\,a_{n+1}\,}{a_n}
    \pause
    = \frac{ar^n}{\,ar^{n-1}\,}
    \pause
    = r
  \]

	\pause
  Se $\rho<1$     \tabto{17mm} a série converge

  \pause
  Se $\rho\geq 1$ \tabto{17mm} a série diverge
\end{proofframe}

%------------------------------------------------------------------------------%
\begin{frame}{Teste da Razão}
	Seja $\;\sum a_n\;$ tal que $\;a_n>0\;$ para $\;n\geq N\;$ \pause e
	\[
     \rho = \lim_{n\to\infty}\frac{a_{n+1}}{a_n}
  \]

  \pause
  \smallskip
	Se $\rho<1$                    \tabto{37mm} a série \emph{converge}

  \pause
  \smallskip
	Se $\rho>1$ ou $\rho\to\infty$ \tabto{37mm} a série \emph{diverge}

  \pause
  \smallskip
	Se $\rho=1$                    \tabto{37mm} o teste é \emph{inconclusivo}
\end{frame}

%------------------------------------------------------------------------------%
\section{Exemplos}

\input{ex/G-06-Teste_razao-ex-1}
\input{ex/G-06-Teste_razao-ex-2}
\input{ex/G-06-Teste_razao-ex-3}
\input{ex/G-06-Teste_razao-ex-4}

%------------------------------------------------------------------------------%
\section{Lista Mínima}

%------------------------------------------------------------------------------%
\begin{frame}{Lista Mínima}
  Estudar as Seção 6.6 da Apostila

  \vspace{\baselineskip}
  Exercícios: 3a-f, 4a-f

  \vfill
  Atenção: A prova é baseada no livro, não nas apresentações
\end{frame}

%------------------------------------------------------------------------------%
\end{document}
%------------------------------------------------------------------------------%
